\documentclass[review,11pt]{article}
\usepackage{geometry}
\geometry{a4paper, margin=1in}
\usepackage{hyperref}
\usepackage{booktabs}
\usepackage{float}
\usepackage{amsmath}
\usepackage{graphicx}
\usepackage{multirow}
\usepackage{xcolor}
\usepackage{amssymb}
\usepackage{subcaption}
\usepackage{caption}

\captionsetup[table]{name=Table S}
\captionsetup[figure]{name=Figure S}

\title{Supplementary Materials: \\ Structure-guided discovery of divergent virulence and resistance candidates in an open pangenome clinical \textit{Escherichia coli} isolate}
\author{Sarra HOSNI-BENMOUMOU and Atika MEKLAT}
\date{}

\begin{document}

\maketitle

\noindent
\textbf{Affiliation:} \\
Laboratoire de Biologie des Systèmes Microbiens, École Normale Supérieure de Kouba, 16000 Alger, Algeria \\
Corresponding author email: \href{mailto:sarra_ben@outlook.fr}{sarra\_ben@outlook.fr}

\vspace{1cm}

\section{Supplementary Tables}

\section{Supplementary Tables}

\begin{table}[H]
\centering
\caption{\textit{De Novo} Draft Genome Assembly and Quality Control Statistics for \textit{E. coli} QA5221 (originally Table 1).}
\label{tab:assembly_qc}
\begin{tabular}{lr}
\toprule
\textbf{Assembly Quality Metric} & \textbf{Value} \\
\midrule
Draft Genome Length & 5,026,316 bp \\
Total Length (Contigs $\ge$ 500 bp) & 5,010,493 bp \\
Total Number of Contigs (All / $\ge$ 500 bp) & 121 / 60 \\
N50 Contig Size & 294,640 bp \\
L50 Contig Count & 7 contigs \\
Largest Assembled Contig & 513,612 bp \\
Assembly GC Content & 50.46\% \\
Genome Fraction Covered (vs. E. coli K-12 MG1655) & 87.06\% \\
Total Unaligned Sequence Length & 969,074 bp \\
Mean Read-Mapping Depth (mosdepth) & 43.89$\times$ \\
\bottomrule
\end{tabular}%
\end{table}

\begin{table}[H]
\centering
\caption{Pangenome Partitioning Metrics of the \textit{E. coli} Expanded Cohort ($n = 400$ Genomes) (originally Table 2).}
\label{tab:pangenome_metrics}
\resizebox{\columnwidth}{!}{%
\begin{tabular}{lccl}
\toprule
\textbf{Pangenome Partition} & \textbf{Strain Frequency} & \textbf{Gene Clusters} & \textbf{Percentage (\%)} \\
\midrule
Core Genome & $\ge$ 99\% (397--401 strains) & 2,650 & 6.14\% \\
Soft Core Genome & 95\%--99\% (381--396 strains) & 640 & 1.48\% \\
Shell Genome & 15\%--95\% (61--380 strains) & 2,448 & 5.68\% \\
Cloud Genome & $< 15\%$ (1--60 strains) & 37,387 & 86.69\% \\
\midrule
\textbf{Total Pangenome} & \textbf{0\%--100\%} & \textbf{43,125} & \textbf{100.00\%} \\
\bottomrule
\end{tabular}%
}
\end{table}

\begin{table}[H]
\centering
\caption{Mobile Genetic Elements (MGEs): IS Families and Integron Cassettes in QA5221 (originally Table 3).}
\label{tab:mge_inventory}
\resizebox{\columnwidth}{!}{%
\begin{tabular}{lcccl}
\toprule
\textbf{MGE Component} & \textbf{Count / Size} & \textbf{bp / Coords} & \textbf{Location} & \textbf{Biological Features} \\
\midrule
\textbf{Insertion Sequences} & 22 elements & 29,217 bp & Multiple contigs & 0.58\% of total genome \\
\quad -- IS3 Family & 5 copies & -- & NODE\_2, 12, 32, 46 & Transposase IS3 family \\
\quad -- ISNCY Family & 4 copies & -- & NODE\_2, 7, 8 & Non-classified elements \\
\quad -- IS200/IS605 Family & 3 copies & $\sim$513 bp & NODE\_7, 27, 60 & Transposase Y1 family \\
\quad -- IS1 Family & 2 copies & 768 bp & NODE\_30 & Flanks active MDR cassettes \\
\midrule
\textbf{Integrons} & & & & \\
\quad -- In0 (Integrase only) & 1 element & 567--1490 & NODE\_43 & Encodes functional \textit{intI1} integrase \\
\quad -- CALIN (Cassette Array) & 1 array & 2--9804 & NODE\_37 & 6 \textit{attC} sites; carries 6 AMR genes \\
\bottomrule
\end{tabular}%
}
\end{table}

\begin{table}[H]
\centering
\caption{Detailed Structural, Genomic, and Mobilome Profiles of the 12 Acquired Resistance Genes in QA5221 (originally Table 5).}
\label{tab:amr_details}
\resizebox{\columnwidth}{!}{%
\begin{tabular}{lllcrrcl}
\toprule
\textbf{Gene Symbol} & \textbf{Locus Tag} & \textbf{AMR Class} & \textbf{Mechanism} & \textbf{Length (bp)} & \textbf{GC (\%)} & \textbf{GC3 (\%)} & \textbf{Location / MGE Context} \\
\midrule
\textit{tet(A)} & KNGPFPPJ\_04581 & Tetracycline & Efflux pump (Class C) & 1200 & 63.67\% & 74.75\% & NODE\_30 (Flanked by IS1; tetR-tetA locus) \\
\textit{estX} & KNGPFPPJ\_04687 & Macrolide & Macrolide esterase & 1080 & 49.17\% & 47.78\% & NODE\_37 (CALIN array cassette 1) \\
\textit{aadA2} & KNGPFPPJ\_04689 & Aminoglycoside & Nucleotidyltransferase & 780 & 51.92\% & 56.15\% & NODE\_37 (CALIN array cassette 2) \\
\textit{cmlA1} & KNGPFPPJ\_04690 & Phenicol & Efflux pump (MFS) & 1260 & 55.32\% & 58.33\% & NODE\_37 (CALIN array cassette 3) \\
\textit{aadA1} & KNGPFPPJ\_04691 & Aminoglycoside & Nucleotidyltransferase & 780 & 52.82\% & 54.23\% & NODE\_37 (CALIN array cassette 4) \\
\textit{qacL} & KNGPFPPJ\_04692 & Biocide & Efflux pump (SMR) & 333 & 51.65\% & 54.95\% & NODE\_37 (CALIN array cassette 5) \\
\textit{sul3} & KNGPFPPJ\_04694 & Sulfonamide & Dihydropteroate synthase & 792 & 37.75\% & 27.27\% & NODE\_37 (CALIN array cassette 6) \\
\textit{dfrA7} & KNGPFPPJ\_04716 & Trimethoprim & Dihydrofolate reductase & 702 & 40.60\% & 37.61\% & NODE\_41 (Class 1 Integron 3'-CS) \\
\textit{qacE$\Delta$1} & KNGPFPPJ\_04717 & Biocide & Efflux pump (SMR) & 348 & 50.00\% & 49.14\% & NODE\_41 (Class 1 Integron 3'-CS distal) \\
\textit{sul1} & KNGPFPPJ\_04718 & Sulfonamide & Dihydropteroate synthase & 780 & 61.92\% & 70.77\% & NODE\_41 (Class 1 Integron 3'-CS distal) \\
\textit{blaTEM-1B} & KNGPFPPJ\_04720 & Beta-lactam & Beta-lactamase (Class A) & 861 & 49.01\% & 43.21\% & NODE\_42 (Tn3-like transposon structure) \\
\textit{aph(3')-Ia} & KNGPFPPJ\_04727 & Aminoglycoside & O-phosphotransferase & 816 & 43.87\% & 40.44\% & NODE\_47 (Isolated resistance element) \\
\bottomrule
\end{tabular}%
}
\end{table}

\begin{table}[H]
\centering
\caption{Resistome Database Synonym Dictionary and Accession Mapping for \textit{E. coli} QA5221 (originally Table 6).}
\label{tab:amr_synonyms}
\resizebox{\columnwidth}{!}{%
\begin{tabular}{lllll}
\toprule
\textbf{Gene Symbol} & \textbf{AMRFinderPlus Name (Acc)} & \textbf{ResFinder Name (Acc)} & \textbf{CARD Name (Acc)} & \textbf{ARG-ANNOT Name (Acc)} \\
\midrule
\textit{tet(A)} & tet(A) (NF012193.0) & tet(A)\_6 (AF534183) & tet(A) (AF534183.1) & (Tet)tetA (JX424423) \\
\textit{estX} & estX (NF033124.1) & -- & -- & -- \\
\textit{aadA2} & aadA2 (NF012157.0) & aadA2\_1 (NC\_010870) & aadA2 (AF156486.1) & (AGly)aadA2 (X68227) \\
\textit{cmlA1} & cmlA1 (NF000509.1) & CmlA7\_1 (JQ639792) & cmlA1 (M64556.1) & (Phe)cmlA1 (M64556) \\
\textit{aadA1} & aadA1 (NF033126.1) & ant(3'')-Ia\_1 (X02340) & aadA (AF550679.1) & (AGly)aadA1-pm (JQ690540) \\
\textit{qacL} & qacL (NF000067.2) & -- & qacL (DQ149925.1) & -- \\
\textit{sul3} & sul3 (NF000296.1) & sul3\_2 (AJ459418) & sul3 (FJ196385.1) & (Sul)sul3 (HQ875016) \\
\textit{dfrA7} & dfrA7 (NF000330.1) & dfrA7\_5 (AJ419170) & dfrA7 (FJ854362.1) & (Tmt)dhfr7 (X58425) \\
\textit{qacE$\Delta$1} & qacEdelta1 (NF000276.2) & -- & qacEdelta1 (U49101.1) & -- \\
\textit{sul1} & sul1 (NA) & sul1\_5 (EU780013) & sul1 (JF969163.1) & (Sul)sul1 (AF071413) \\
\textit{blaTEM-1B} & blaTEM-1 (NF000531.2) & blaTEM-1B\_1 (AY458016) & TEM-1 (AL513383.1) & (Bla)blaTEM-105 (AF516720) \\
\textit{aph(3')-Ia} & aph(3')-Ia (NF051474.1) & aph(3')-Ia\_1 (V00359) & APH(3')-Ia (BX664015.1) & (AGly)aph3-Ia (HQ840942) \\
\bottomrule
\end{tabular}%
}
\end{table}

\begin{table}[H]
\centering
\caption{Statistical Comparisons of Gene GC\%, GC3\% Codon Bias, and Length across Partitions (originally Table 7).}
\label{tab:stats_comparisons}
\resizebox{\columnwidth}{!}{%
\begin{tabular}{lcccc}
\toprule
\textbf{Metric / Feature} & \textbf{Core ($N=3,310$)} & \textbf{Accessory ($N=1,047$)} & \textbf{Singleton ($N=278$)} & \textbf{Acquired AMR ($N=12$)} \\
\midrule
\textbf{GC Content (\%)} & $51.86\% \pm 3.72\%$ & $48.49\% \pm 6.90\%$ & $46.78\% \pm 7.86\%$ & $50.64\% \pm 7.65\%$ \\
\quad -- Shapiro-Wilk $p$-value & $2.387e-17$ & $1.116e-15$ & $1.119e-04$ & $7.873e-01$ \\
\quad -- Mann-Whitney $p$ vs. Core & -- & $7.812e-44$ & $8.431e-31$ & $1.000e+00$ \\
\quad -- Mann-Whitney $p$ vs. AMR & $1.000e+00$ & $1.000e+00$ & $9.360e-01$ & -- \\
\midrule
\textbf{GC3 Codon Bias (\%)} & $55.90\% \pm 6.63\%$ & $50.36\% \pm 11.33\%$ & $48.03\% \pm 12.89\%$ & $51.22\% \pm 13.45\%$ \\
\quad -- Shapiro-Wilk $p$-value & $9.872e-10$ & $8.801e-06$ & $7.313e-02$ & $9.705e-01$ \\
\quad -- Mann-Whitney $p$ vs. Core & -- & $6.269e-50$ & $1.462e-28$ & $4.224e-01$ \\
\quad -- Mann-Whitney $p$ vs. AMR & $4.224e-01$ & $1.000e+00$ & $1.000e+00$ & -- \\
\midrule
\textbf{Gene Length (bp)} & $981.9 \pm 626.7$ & $940.4 \pm 797.6$ & $677.0 \pm 561.1$ & $811.0 \pm 283.8$ \\
\quad -- Shapiro-Wilk $p$-value & $2.824e-25$ & $2.185e-37$ & $7.587e-18$ & $1.919e-01$ \\
\quad -- Mann-Whitney $p$ vs. Core & -- & $8.179e-05$ & $7.368e-23$ & $1.000e+00$ \\
\quad -- Mann-Whitney $p$ vs. AMR & $1.000e+00$ & $1.000e+00$ & $2.726e-01$ & -- \\
\bottomrule
\end{tabular}%
}
\end{table}

% Table S7: MM-GBSA Extended Parameters
\begin{table}[H]
\centering
\caption{MM-GBSA Binding Free Energy Calculation: Extended Parameter Settings and Energy Decomposition for the GNAT Candidate \textit{KNGPFPPJ\_02769}.}
\label{tab:mmgbsa_params}
\begin{tabular}{ll}
\toprule
\textbf{Parameter} & \textbf{Value / Description} \\
\midrule
\multicolumn{2}{l}{\textit{Solvation Model}} \\
\quad Model & Still Generalized Born (GB) implicit solvent \\
\quad Solvent dielectric ($\varepsilon_{\text{solv}}$) & 80 (water at 298 K) \\
\quad Solute dielectric ($\varepsilon_{\text{solute}}$) & 1 \\
\midrule
\multicolumn{2}{l}{\textit{Born Radii (per element)}} \\
\quad H (hydrogen) & 1.20 \AA \\
\quad C (carbon) & 1.70 \AA \\
\quad N (nitrogen) & 1.55 \AA \\
\quad O (oxygen) & 1.52 \AA \\
\quad S (sulfur) & 1.80 \AA \\
\midrule
\multicolumn{2}{l}{\textit{van der Waals Interactions}} \\
\quad Potential & Lennard-Jones 12-6 \\
\quad Per-pair energy cap & $+$2 kcal/mol (clash correction) \\
\midrule
\multicolumn{2}{l}{\textit{Nonpolar Solvation}} \\
\quad Method & Solvent-Accessible Surface Area (SASA) \\
\quad Surface tension coefficient ($\gamma$) & $-0.0054$ kcal/mol/\AA$^2$ \\
\midrule
\multicolumn{2}{l}{\textit{Charges}} \\
\quad Receptor protein & AMBER-parameterized backbone partial charges \\
\quad Ligand & AutoDock Vina Gasteiger charges (from docked PDBQT) \\
\midrule
\multicolumn{2}{l}{\textit{Energy Components}} \\
\quad $\Delta E_{\text{elec}}$ & Coulomb electrostatic interaction energy \\
\quad $\Delta E_{\text{vdW}}$ & Lennard-Jones 12-6 van der Waals energy \\
\quad $\Delta G_{\text{pol}}$ & Polar (Still GB) solvation free energy change \\
\quad $\Delta G_{\text{npol}}$ & Nonpolar solvation free energy (SASA-based) \\
\quad $\Delta G_{\text{bind}}$ & Total binding free energy (sum of above) \\
\bottomrule
\end{tabular}
\end{table}

% Table S8: BLASTP Homologs
\begin{table}[H]
\centering
\caption{Top 20 BLASTP Homologs of \textit{KNGPFPPJ\_02769} in the NCBI nr Database ($n = 200$ total significant hits; E $< 10^{-5}$, BLOSUM62 matrix). Hits are ranked by sequence identity. All top hits originate from Proteobacteria (Enterobacteriaceae), with 100\% identity to \textit{E.~coli} and \textit{Salmonella enterica} strains.}
\label{tab:blast_top20}
\resizebox{\columnwidth}{!}{%
\begin{tabular}{rlllrr}
\toprule
\textbf{Rank} & \textbf{Organism} & \textbf{Phylum} & \textbf{Identity (\%)} & \textbf{E-value} & \textbf{Coverage (\%)} \\
\midrule
1  & \textit{Escherichia coli}    & Proteobacteria & 100.0 & 0.00      & 99.7 \\
2  & \textit{Escherichia coli}    & Proteobacteria & 100.0 & 0.00      & 81.2 \\
3  & \textit{Escherichia coli}    & Proteobacteria & 100.0 & 3.27e-149 & 58.7 \\
4  & \textit{Escherichia coli}    & Proteobacteria & 100.0 & 2.03e-142 & 56.1 \\
5  & \textit{Escherichia coli}    & Proteobacteria & 100.0 & 7.91e-141 & 55.6 \\
6  & \textit{Escherichia coli}    & Proteobacteria & 100.0 & 5.91e-140 & 55.3 \\
7  & \textit{Salmonella enterica} & Proteobacteria & 100.0 & 3.88e-139 & 55.0 \\
8  & \textit{Escherichia coli}    & Proteobacteria & 100.0 & 8.32e-130 & 51.6 \\
9  & \textit{Salmonella enterica} & Proteobacteria & 100.0 & 1.44e-126 & 50.4 \\
10 & \textit{Escherichia coli}    & Proteobacteria & 100.0 & 3.81e-121 & 48.4 \\
11 & \textit{Escherichia coli}    & Proteobacteria & 99.7  & 0.00      & 99.7 \\
12 & \textit{Escherichia coli}    & Proteobacteria & 99.7  & 0.00      & 99.7 \\
13 & \textit{Escherichia coli}    & Proteobacteria & 99.7  & 0.00      & 95.2 \\
14 & \textit{Escherichia coli}    & Proteobacteria & 99.6  & 0.00      & 75.5 \\
15 & \textit{Escherichia coli}    & Proteobacteria & 99.5  & 6.75e-156 & 61.3 \\
16 & \textit{Escherichia coli}    & Proteobacteria & 99.5  & 2.35e-150 & 59.3 \\
17 & \textit{Escherichia coli}    & Proteobacteria & 99.5  & 5.38e-148 & 58.4 \\
18 & \textit{Escherichia coli}    & Proteobacteria & 99.5  & 5.26e-146 & 57.8 \\
19 & \textit{Escherichia coli}    & Proteobacteria & 99.5  & 5.77e-144 & 57.0 \\
20 & \textit{Escherichia coli}    & Proteobacteria & 99.5  & 4.28e-142 & 56.1 \\
\bottomrule
\multicolumn{6}{l}{\footnotesize Full hit table available on GitHub: \href{https://github.com/hosniadilemp-a11y/AMR_Novel_Gene_Discovery/blob/main/results/Step7/KNGPFPPJ_02769_blast_hits.tsv}{\texttt{results/Step7/KNGPFPPJ\_02769\_blast\_hits.tsv}}} \\
\end{tabular}%
}
\end{table}

\clearpage
% Table S9: Reconciliation Matrix
\begin{table*}[p]
\centering
\caption{Reconciliation Matrix of Functional/Structural Inclusion and Filtering Thresholds for all 22 Prioritized Candidates.}
\label{tab:candidate_reconciliation}
\resizebox{\textwidth}{!}{%
\begin{tabular}{lcccccccp{7cm}}
\toprule
\textbf{Locus Tag} & \textbf{Size (aa)} & \textbf{Length Filter} & \textbf{Pfam Hits} & \textbf{GC Content} & \textbf{GC Depression} & \textbf{TM-Score} & \textbf{Structural Homology} & \textbf{MGE / Genomic Island Context} \\
\midrule
KNGPFPPJ\_02769 & 351 & Pass & 0 (Pass) & 32.9\% & Pass ($p < 10^{-4}$) & 0.9569 (Pass) & Pass & NODE\_9 Composite Transposon \\
KNGPFPPJ\_00061 & 350 & Pass & 0 (Pass) & 48.5\% & Pass ($p < 10^{-4}$) & 0.5899 (Pass) & Pass & NODE\_1 Intact Prophage \\
KNGPFPPJ\_03161 & 372 & Pass & 0 (Pass) & 29.3\% & Pass ($p < 10^{-4}$) & 0.8953 (Pass) & Pass & NODE\_1 Intact Prophage \\
KNGPFPPJ\_04371 & 984 & Pass & 0 (Pass) & 35.8\% & Pass ($p < 10^{-4}$) & 0.3450 (Domain Pass: 0.72) & Pass & NODE\_24 Plasmid mobilization array \\
KNGPFPPJ\_00084 & 203 & Pass & 0 (Pass) & 45.8\% & Pass ($p < 10^{-4}$) & 0.5352 (Pass) & Pass & NODE\_1 Intact Prophage \\
KNGPFPPJ\_00091 & 249 & Pass & 0 (Pass) & 51.2\% & Pass ($p < 0.05$) & 0.4573 (Partial Pass) & Pass & NODE\_1 Intact Prophage \\
KNGPFPPJ\_00107 & 216 & Pass & 0 (Pass) & 48.5\% & Pass ($p < 10^{-4}$) & 0.5864 (Pass) & Pass & NODE\_1 Intact Prophage \\
KNGPFPPJ\_00109 & 344 & Pass & 0 (Pass) & 44.9\% & Pass ($p < 10^{-4}$) & 0.7750 (Pass) & Pass & NODE\_1 Intact Prophage \\
KNGPFPPJ\_00114 & 399 & Pass & 0 (Pass) & 48.6\% & Pass ($p < 10^{-4}$) & 0.4854 (Partial Pass) & Pass & NODE\_1 Intact Prophage \\
KNGPFPPJ\_04325 & 294 & Pass & 0 (Pass) & 35.0\% & Pass ($p < 10^{-4}$) & 0.5000 (Pass) & Pass & NODE\_23 Recombinase-associated Island \\
KNGPFPPJ\_04326 & 501 & Pass & 0 (Pass) & 35.1\% & Pass ($p < 10^{-4}$) & 0.5200 (Pass) & Pass & NODE\_23 Recombinase-associated Island \\
KNGPFPPJ\_03157 & 371 & Pass & 0 (Pass) & 31.4\% & Pass ($p < 10^{-4}$) & 0.8972 (Pass) & Pass & NODE\_1 Intact Prophage \\
KNGPFPPJ\_03836 & 246 & Pass & 0 (Pass) & 46.7\% & Pass ($p < 10^{-4}$) & 0.9502 (Pass) & Pass & NODE\_1 Intact Prophage \\
KNGPFPPJ\_00097 & 313 & Pass & 0 (Pass) & 53.1\% & Pass ($p < 0.05$) & 0.8229 (Pass) & Pass & NODE\_1 Intact Prophage \\
KNGPFPPJ\_00103 & 381 & Pass & 0 (Pass) & 49.0\% & Pass ($p < 10^{-4}$) & 0.9051 (Pass) & Pass & NODE\_1 Intact Prophage \\
KNGPFPPJ\_00106 & 668 & Pass & 0 (Pass) & 52.2\% & Pass ($p < 0.05$) & 0.3253 (Partial Pass) & Pass & NODE\_1 Intact Prophage \\
KNGPFPPJ\_00112 & 250 & Pass & 0 (Pass) & 51.1\% & Pass ($p < 0.05$) & 0.8974 (Pass) & Pass & NODE\_1 Intact Prophage \\
KNGPFPPJ\_00115 & 226 & Pass & 0 (Pass) & 47.9\% & Pass ($p < 10^{-4}$) & 0.6774 (Pass) & Pass & NODE\_1 Intact Prophage \\
KNGPFPPJ\_01156 & 300 & Pass & 0 (Pass) & 32.6\% & Pass ($p < 10^{-4}$) & 0.7379 (Pass) & Pass & Accessory chromosomal island \\
KNGPFPPJ\_01158 & 677 & Pass & 0 (Pass) & 32.6\% & Pass ($p < 10^{-4}$) & 0.7823 (Pass) & Pass & Accessory chromosomal island \\
KNGPFPPJ\_04303 & 256 & Pass & 0 (Pass) & 39.4\% & Pass ($p < 10^{-4}$) & 0.7623 (Pass) & Pass & NODE\_23 Recombinase-associated Island \\
KNGPFPPJ\_04571 & 290 & Pass & 0 (Pass) & 38.6\% & Pass ($p < 10^{-4}$) & 0.7169 (Pass) & Pass & Accessory chromosomal island \\
\bottomrule
\end{tabular}%
}
\end{table*}

% Table S7: Extended MM-GBSA Parameters & Alanine Scanning
\begin{table}[H]
\centering
\caption{\textbf{MM-GBSA Binding Free Energy Decomposition and \textit{In Silico} Active-Site Alanine Scanning for GNAT\_KA27 Bound to Kanamycin.}}
\label{tab:alanine_scanning}
\resizebox{\columnwidth}{!}{%
\begin{tabular}{lccccc}
\toprule
\textbf{Variant} & \textbf{Receptor $\Delta G$ (kcal/mol)} & \textbf{Ligand $\Delta G$ (kcal/mol)} & \textbf{Complex $\Delta G$ (kcal/mol)} & \textbf{Total $\Delta G_{\text{bind}}$ (kcal/mol)} & \textbf{$\Delta\Delta G_{\text{mut}}$ (kcal/mol)} \\
\midrule
Wild-Type (WT)     & $-14250.30$ & $-45.20$ & $-14319.40$ & $-23.90$ & $0.00$ \\
L125A              & $-14248.10$ & $-45.20$ & $-14314.30$ & $-21.00$ & $+2.90$ \\
L187A              & $-14247.30$ & $-45.20$ & $-14312.73$ & $-20.23$ & $+3.67$ \\
F188A              & $-14245.80$ & $-45.20$ & $-14309.65$ & $-18.65$ & $+5.25$ \\
Triple Mutant      & $-14239.10$ & $-45.20$ & $-14296.38$ & $-12.08$ & $+11.82$ \\
\bottomrule
\end{tabular}%
}
\end{table}

\begin{table}[H]
\centering
\caption{\textbf{Consensus Structural Homology Alignment Metrics.}}
\label{tab:consensus_structural_alignment}
\resizebox{\columnwidth}{!}{%
\begin{tabular}{llccc}
\toprule
\textbf{Candidate} & \textbf{Top Homolog Template} & \textbf{Foldseek TM-Score} & \textbf{CEalign RMSD (\AA)} & \textbf{PDB Match} \\
\midrule
GNAT\_KA27 & \textit{S. enterica} GNAT & 0.9569 & 1.53 & AF-A0A3U7PQ76-F1 \\
Ehly\_61   & Putative Enterohemolysin & 0.5900 & 4.27 & AF-A0A2Y8JZY2-F1 \\
OAgP\_161  & Wzy Polymerase           & 0.8953 & 2.53 & AF-A0A743B7N1-F1 \\
OAT\_371   & Autotransporter Barrel   & 0.7200 & 3.03 & AF-P33924-F1 \\
\bottomrule
\end{tabular}%
}
\end{table}

\clearpage
\section{Supplementary Figures}

\begin{figure}[H]
\centering
\includegraphics[width=0.75\textwidth]{plots/Fig02_Assembly_Coverage.pdf}
\caption{Contig Length versus Read Mapping Coverage Scatter Plot for \textit{E. coli} QA5221 (originally Figure 2).}
\label{fig:assembly_qc}
\end{figure}

\begin{figure}[H]
\centering
\includegraphics[width=0.85\textwidth]{plots/Fig03-c_Pangenome_Curves.pdf}
\caption{Heap's Law pangenome accumulation and new gene discovery curves for the \textit{E. coli} ST354 lineage ($n = 32$) (originally Figure 3). (a) Pangenome accumulation curve showing total gene clusters as a function of the number of genomes sampled. (b) New genes discovery curve showing the mean number of new genes identified per genome added.}
\label{fig:pangenome_curves}
\end{figure}

\begin{figure}[H]
\centering
\includegraphics[width=0.85\textwidth]{plots/Fig03-b_Gene_Frequency_Stats.png}
\caption{Cohort Prevalence and Count Distribution of the four prioritized candidates across the expanded 400-genome cohort (originally Figure 3-b). (a) Percentage occurrence showing GNAT and Wzy as private singletons, and hemolysin and autotransporter as rare accessory genes. (b) Genome count distribution between the clinical isolate QA5221 (red) and the comparative public genomes (blue).}
\label{fig:gene_frequencies}
\end{figure}

\begin{figure}[H]
\centering
\includegraphics[width=0.75\textwidth]{plots/Fig04-b_Virulence_Heatmap.pdf}
\caption{Heatmap demonstrating the distribution of variable virulence-associated genes (top 40 variable genes) across the ST354 lineage (originally Figure 4-b). Node colors represent presence (1, dark blue) or absence (0, light yellow) of each determinant.}
\label{fig:virulence_heatmap}
\end{figure}

\begin{figure}[H]
\centering
\includegraphics[width=0.65\textwidth]{plots/Fig05_IS_Distribution.pdf}
\caption{Insertion Sequence (IS) Families Abundance in \textit{E. coli} QA5221 (originally Figure 5).}
\label{fig:is_distribution}
\end{figure}

\begin{figure}[H]
\centering
\includegraphics[width=0.75\textwidth]{plots/Fig06-b_AMR_Network.pdf}
\caption{Antimicrobial resistance gene co-occurrence network across the ST354 clonal complex (originally Figure 6-b). Node sizes represent carriage frequency in the cohort, and edge weights represent Jaccard similarity.}
\label{fig:amr_network}
\end{figure}

\begin{figure}[H]
\centering
\includegraphics[width=0.75\textwidth]{plots/Fig07_AMR_Comparison.pdf}
\caption{Comparative Boxplots of Genomic Features with Acquired AMR Genes Included (originally Figure 7).}
\label{fig:amr_comparison}
\end{figure}

\begin{figure}[H]
\centering
\includegraphics[width=0.75\textwidth]{plots/Fig12_GNAT_Phylogenetic_Tree.pdf}
\caption{Maximum-Likelihood phylogenetic tree of the GNAT/AAC superfamily showing KNGPFPPJ\_02769 ($\bigstar$) branching as a sister clade to clinical AACs (originally Figure 12).}
\label{fig:gnat_tree}
\end{figure}

\begin{figure}[H]
\centering
\includegraphics[width=0.75\textwidth]{plots/Fig15-a_Hemolysin_Phylogenetic_Tree.pdf}
\caption{Maximum-Likelihood family-level phylogenetic tree showing KNGPFPPJ\_00061 ($\bigstar$) branching as an isolated sister clade within the hemolysin/pore-forming toxin superfamily (originally Figure 15-a).}
\label{fig:hemolysin_phylo_tree}
\end{figure}

\begin{figure}[H]
\centering
\includegraphics[width=0.75\textwidth]{plots/Fig15-b_Wzy_Phylogenetic_Tree.pdf}
\caption{Maximum-Likelihood family-level phylogenetic tree showing KNGPFPPJ\_03161 ($\bigstar$) branching separately from canonical \textit{E. coli} Wzy polymerases (originally Figure 15-b).}
\label{fig:wzy_phylo_tree}
\end{figure}

\begin{figure}[H]
\centering
\includegraphics[width=0.75\textwidth]{plots/Fig15-c_Autotransporter_Phylogenetic_Tree.pdf}
\caption{Maximum-Likelihood family-level phylogenetic tree showing KNGPFPPJ\_04371 ($\bigstar$) as an isolated outgroup relative to canonical adhesin and SPATE subfamilies (originally Figure 15-c).}
\label{fig:autotr_phylo_tree}
\end{figure}

\begin{figure}[H]
\centering
\includegraphics[width=0.75\textwidth]{plots/Fig18-a_Structural_Alignment_RMSD.pdf}
\caption{Comparative structural alignment Root Mean Square Deviation (RMSD) for the four prioritised candidates versus their templates (originally Figure 18-a). Dashed line at 3.0~\AA\ marks the acceptable structural fit threshold.}
\label{fig:struct_rmsd}
\end{figure}

\begin{figure}[H]
\centering
\includegraphics[width=0.75\textwidth]{plots/Fig18-b_Structural_Alignment_Twilight_Zone.pdf}
\caption{Scatter plot of sequence identity versus TM-score illustrating the ``twilight zone of homology'' (originally Figure 18-b): candidates below 25\% sequence identity but above TM-score 0.50 are detectable only by structure-based methods.}
\label{fig:struct_twilight_zone}
\end{figure}

\begin{figure}[H]
\centering
\includegraphics[width=0.75\textwidth]{plots/Fig20-b_PLM_Candidate_Highlight_t-sne_projection.pdf}
\caption{ESM-2 650M Protein Language Model (\texttt{esm2\_t33\_650M\_UR50D}, 1280-dim) t-SNE projection of candidate and reference proteins. Stars ($\bigstar$) denote the four novel candidates (red = GNAT; green = hemolysin; yellow = Wzy; blue = autotransporter). Diamond markers indicate acquired AMR reference genes. In the t-SNE embedding projection, the GNAT candidate KNGPFPPJ\_02769 co-localizes with known aminoglycoside resistance determinants, confirming the robust topological grouping observed in the main UMAP projection.}
\label{fig:plm_tsne_supp}
\end{figure}

\begin{figure}[H]
\centering
\includegraphics[width=0.85\textwidth]{../figures/FigS24_DSSP_Timeline.png}
\caption{\textbf{GNAT\_KA27 Explicit-Solvent MD DSSP Secondary Structure Timeline.} Secondary structure content ($\alpha$-helix, $\beta$-sheet, coil/loop) calculated over the 127.78~ns trajectory. Mean values: $29.1\%$ $\alpha$-helix, $30.5\%$ $\beta$-sheet, and $40.4\%$ coil/loop, confirming fold stability without unfolding.}
\label{fig:dssp_timeline}
\end{figure}

\begin{figure}[H]
\centering
\includegraphics[width=0.95\textwidth]{../figures/FigS25_PAE_Heatmaps.png}
\caption{\textbf{AlphaFold3 and ESMFold 2D Predicted Aligned Error (PAE) and pLDDT Confidence Heatmaps.} (a) GNAT\_KA27 (mean pLDDT 87.2\%), (b) Ehly\_61 (mean pLDDT 88.4\%), (c) OAgP\_161 (mean pLDDT 81.5\%), and (d) OAT\_371 (mean pLDDT 84.2\%; bimodal profile showing high confidence in the C-terminal $\beta$-barrel and higher uncertainty in the extracellular passenger domain).}
\label{fig:pae_heatmaps}
\end{figure}

\begin{figure}[H]
\centering
\includegraphics[width=0.85\textwidth]{../figures/FigS26_Foldseek_False_Positive_Control.png}
\caption{\textbf{Empirical Foldseek False-Positive Control Rate Distribution.} Distribution of maximum TM-scores for 10 random non-homologous accessory singletons (grey circles) compared to the four prioritized candidates (red triangles). All 10 control singletons fall below the empirical noise limit of $\text{TM} = 0.40$, validating the $\text{TM} \ge 0.50$ prioritization cutoff.}
\label{fig:false_positive_control}
\end{figure}

\begin{figure}[H]
\centering
\includegraphics[width=0.95\textwidth]{../figures/FigS27_Alanine_Scanning_Energetics.png}
\caption{\textbf{\textit{In Silico} Active-Site Alanine Scanning Binding Energetics for GNAT\_KA27 Bound to Kanamycin.} (a) MM-GBSA total binding free energy ($\Delta G_{\text{bind}}$) across wild-type (WT) and mutant complexes relative to non-specific decoy baseline ($-12.14$~kcal/mol). (b) Mutation affinity loss ($\Delta\Delta G_{\text{mut}}$), showing progressive loss of binding free energy upon mutating Leu125, Leu187, and Phe188 to alanine.}
\label{fig:alanine_scanning_plot}
\end{figure}

\begin{figure}[H]
\centering
\includegraphics[width=0.98\textwidth]{../figures/FigS28_Replicate_MD_RMSD.png}
\caption{\textbf{Multi-Candidate Replicate Explicit-Solvent MD Simulations ($n=3$ Independent Replicates Each).} Backbone C$\alpha$ RMSD trajectories across three independent simulation seeds (blue, green, orange) and ensemble mean trajectories (red bold line) with Block SEM error boundaries for: (a) Soluble acetyltransferase GNAT\_KA27 ($\text{Mean RMSD} = 2.25 \pm 0.06$~\AA, Block SEM $= 0.007$~\AA), (b) Membrane-embedded enterohemolysin Ehly\_61 ($\text{Mean RMSD} = 5.81 \pm 0.22$~\AA, Block SEM $= 0.012$~\AA), and (c) Multi-transmembrane O-antigen polymerase OAgP\_161 ($\text{Mean RMSD} = 2.73 \pm 0.14$~\AA, Block SEM $= 0.015$~\AA). Low-alpha background lines represent raw high-frequency fluctuations.}
\label{fig:replicate_md_rmsd}
\end{figure}

\begin{figure}[H]
\centering
\includegraphics[width=0.85\textwidth]{../figures/FigS29_Comparative_MMGBSA_Canonical_AACs.png}
\caption{\textbf{Comparative MM-GBSA Binding Free Energies of GNAT\_KA27 versus Canonical Resistance Enzymes.} MM-GBSA total binding free energy ($\Delta G_{\text{bind}}$) for GNAT\_KA27 ($-23.90 \pm 1.10$~kcal/mol) compared to canonical AAC(3)-Ib ($-24.52 \pm 0.85$~kcal/mol), APH(3')-Ia ($-26.15 \pm 0.92$~kcal/mol), and non-specific decoy control ($-12.14 \pm 1.69$~kcal/mol).}
\label{fig:comparative_mmgbsa_canonical}
\end{figure}

\begin{figure}[H]
\centering
\includegraphics[width=0.85\textwidth]{../figures/FigS30_OAT371_Barrel_MD_RMSD.png}
\caption{\textbf{OAT\_371 Translocation $\beta$-Barrel Domain Explicit-Solvent MD Dynamics (20~ns).} C$\alpha$ RMSD of the isolated 250-aa C-terminal 12-stranded $\beta$-barrel domain (residues 734--984), showing structural rigidity ($\text{RMSD} = 1.53 \pm 0.09$~\AA) and outer-membrane pore stability.}
\label{fig:oat371_barrel_md}
\end{figure}

\begin{figure}[H]
\centering
\includegraphics[width=0.85\textwidth]{../figures/FigS31_Ensemble_Docking_Energetics.png}
\caption{\textbf{GNAT\_KA27 Ensemble Docking Binding Affinity Across 10 Dynamic Trajectory Snapshots.} Binding free energy calculated across 10 snapshot conformations extracted from 0 to 120~ns ($\text{Mean } \Delta G_{\text{bind}} = -23.08 \pm 0.56$~kcal/mol), proving binding pocket flexibility and persistent high-affinity substrate engagement.}
\label{fig:ensemble_docking_energetics}
\end{figure}

\clearpage
\section{Bioinformatics Software and Python Packages}

% Table S10: Command-line bioinformatics tools
\begin{table*}[h!]
\centering
\caption{Command-line bioinformatics tools used in this study (originally Table 9).}
\label{tab:bioinformatics_tools}
\footnotesize
\begin{tabular}{llp{11cm}}
\toprule
\textbf{Software Tool} & \textbf{Version} & \textbf{Role in Pipeline} \\
\midrule
\href{https://www.bioinformatics.babraham.ac.uk/projects/fastqc/}{FastQC} & v0.12.1 & Read quality control \\
\href{https://cutadapt.readthedocs.io/}{Cutadapt} & v5.2 & Adapter and quality trimming \\
\href{https://multiqc.info/}{MultiQC} & v1.35 & Aggregating quality control reports \\
\href{https://cab.spbu.ru/software/spades/}{SPAdes} & v3.15.5 & \textit{De novo} genome assembly \\
\href{https://quast.sourceforge.net/}{QUAST} & v5.2 & Assembly quality assessment \\
\href{https://github.com/lh3/minimap2}{minimap2} & v2.24 & Read mapping to assembly \\
\href{https://www.htslib.org/}{samtools} & v1.16 & Alignment file processing \\
\href{https://github.com/brentp/mosdepth}{mosdepth} & v0.3.3 & Coverage depth calculation \\
\href{https://github.com/tseemann/prokka}{Prokka} & v1.14.6 & Genome annotation \\
\href{https://github.com/tseemann/abricate}{ABRicate} & v1.0.1 & Contig screening for AMR/virulence \\
\href{https://github.com/ncbi/amr}{AMRFinderPlus} & v4.2.7 & Resistance gene and mutation detection \\
\href{https://github.com/phac-nml/mob-suite}{MOB-recon} & v1.4.9 & Plasmid typing and reconstruction \\
\href{https://gtonkinhill.github.io/panaroo/}{Panaroo} & v1.3.4 & Pangenome orthology clustering \\
\href{http://www.iqtree.org/}{IQ-TREE} & v2.2.0 & Maximum-likelihood phylogenomics \\
\href{https://mafft.cbrc.jp/alignment/software/}{MAFFT} & v7.520 & Multiple sequence alignment \\
\href{https://github.com/xiezhq/ISEScan}{ISEScan} & v1.7.2 & Insertion sequence detection \\
\href{https://integronfinder.readthedocs.io/}{IntegronFinder} & v2.0.2 & Integron identification \\
\href{https://github.com/gamcil/clinker}{clinker} & v0.0.32 & Gene cluster synteny visualization \\
\bottomrule
\end{tabular}
\end{table*}

% Table S11: Bioinformatics and statistical Python packages
\begin{table*}[h!]
\centering
\caption{Bioinformatics and statistical Python packages used in this study (originally Table 10).}
\label{tab:python_packages}
\footnotesize
\begin{tabular}{llp{11cm}}
\toprule
\textbf{Python Package} & \textbf{Version} & \textbf{Role in Pipeline} \\
\midrule
\href{https://biopython.org/}{BioPython} & v1.81 & Sequence manipulation and parsing \\
\href{https://numpy.org/}{NumPy} & v1.24.3 & Numerical array calculations \\
\href{https://scipy.org/}{SciPy} & v1.10.1 & Statistical tests and distribution analysis \\
\href{https://pandas.pydata.org/}{pandas} & v2.0.3 & Data manipulation and tabular analysis \\
\href{https://matplotlib.org/}{matplotlib} & v3.7.1 & Publication-grade plotting \\
\href{https://seaborn.pydata.org/}{seaborn} & v0.12.2 & Statistical data visualization \\
\href{https://www.mdanalysis.org/}{MDAnalysis} & v2.5.0 & Molecular dynamics trajectory parsing \\
\href{https://networkx.org/}{networkx} & v3.1 & Network and graph algorithms \\
\href{https://www.statsmodels.org/}{statsmodels} & v0.14.0 & Statistical models and regressions \\
\href{https://openmm.org/}{openmm} & v8.1.1 & GPU-accelerated molecular dynamics \\
\href{https://github.com/openmm/pdbfixer}{pdbfixer} & v1.9 & PDB file preparation and cleaning \\
\href{https://scikit-learn.org/}{scikit-learn} & v1.2.2 & Dimensionality reduction and clustering \\
\href{https://github.com/lmcinnes/umap}{umap-learn} & v0.5.3 & UMAP projections of embeddings \\
\bottomrule
\end{tabular}
\end{table*}

\end{document}